\documentclass{jnm}
\usepackage[utf8]{inputenc}
\usepackage[T1]{fontenc}
\usepackage{geometry}
\usepackage{graphicx}
\usepackage{float}
\usepackage{hyperref}
\usepackage[english,indonesian]{babel}

\title{Optimasi Pemrosesan Template LaTeX Menggunakan Pendekatan Berbasis AI}
\author{Academic Support Agent}
\date{\today}

\begin{document}
\maketitle

\begin{abstract}

\end{abstract}

\chapter{PENDAHULUAN}


\chapter{TINJAUAN PUSTAKA}


% Original content (commented out):
% %%%%%%%%%%%%%%%%%%%%%%%%%%%%%%%%%%%%%%%%%%%%%%%%%%
%-------------------------------------------------
%	TEMPLATE JNM 2022 - LIMOGES
%	d'après le template JNM2017 de Kevin NADAUD
%-------------------------------------------------
%%%%%%%%%%%%%%%%%%%%%%%%%%%%%%%%%%%%%%%%%%%%%%%%%%

\documentclass{jnm}
\usepackage[T1]{fontenc}
\usepackage[utf8]{inputenc}
\usepackage[cmex10]{amsmath}
\usepackage{array}
\usepackage{graphicx}
\usepackage{caption}
\usepackage{subcaption}
\usepackage{hyperref}
\graphicspath{{./Figures/}}


%---------------------------------------------
% TITRE DE L'ARTICLE
%---------------------------------------------
\title{Titre de l'article}
\author{A. Auteur\textsuperscript{1}, B.-C. Co-auteur\textsuperscript{2}}
\email{auteur\_referent@institution.fr}
\institute{\textsuperscript{1}Laboratoire XLIM, 123 avenue Albert Thomas, 87060 \\\textsuperscript{2}Université de Limoges, 33 rue François Mitterrand, 87032 Limoges\\}


\begin{document}
\maketitle

%---[RESUME]------------------------------------------
\begin{abstract}
Afin d'uniformiser la présentation des papiers, nous vous proposons un guide tant pour la soumission initiale de 2 pages que pour le papier final de 4 pages (si acceptation). Les communications seront soumises uniquement sous format PDF. Il est souhaitable que le résumé soit identique à celui qui sera saisi en ligne lors de la soumission de la version finale de votre papier (max strict 15 lignes). Un soin particulier est demandé aux auteurs pour ce résumé car ce texte sera distribué aux participants en même temps que le programme.
\end{abstract}


%---[1. INTRODCTION] ------------------------------------------
\section{Introduction}
Le format du papier est le format A4. Les marges droite et gauche sont de 2 cm, la marge haute de 2,5 cm et la marge basse de 2 cm. L'en-tête est de 1,25 cm tout comme pour le pied-de-page, SAUF pour la première page où l'en-tête est de 2,5 cm, \og titre de la conférence et logos obligent \fg. Le Tableau~\ref{tab:marges} résume ces caractéristiques.

Titre, auteur et affiliation sont centrés sur la page. Le titre de la communication est en police Times 16, gras. Auteurs et affiliation sont en Times 10, italique. Pour les espacement et enchaînements, vous pouvez vous référer à chaque fois au critère « paragraphe » du style utilisé. Le corps du papier est en deux colonnes de 8,12 cm chacune, espacées de 0,75 cm. Le corps de texte est en police Times 10, avec un retrait de 0,5 cm. Les titres des paragraphes sont en gras, police Times 12. Les titres sont numérotés, avec un retrait de 1 cm. Les sous-titres sont en italique, Times 10 avec deux retraits de 0,5 et 1,5 cm. En utilisant ce modèle Latex, vous n'avez pas à vous soucier de ces différents éléments de mise en forme.

Le paragraphe d'introduction doit traiter du contexte du sujet dans sa globalité avant d'indiquer le schéma précis dans lequel s'inscrit l'article. Ce format d'article s'appuie de manière importante sur le format utilisé lors des précédentes JNM et en particulier celle de 2015 à Bordeaux \cite{kerherve}.



%---[2. Figures, tableaux, légendes et équations]--------------- 
\section{Figures, tableaux, légendes et équations}
\subsection{Figures}

Il est souvent pratique d'insérer les figures flottantes pour laisser le positionnement libre, tel que sur la Figure~\ref{fig:photo}. Cela évite les blancs à la fins des colonnes.

\begin{figure}[t]
  \includegraphics[width=\linewidth]{Limoges.png}
   \caption{Centre de la conférence à Limoges.}
  \label{fig:photo}
\end{figure}
\begin{figure}[t]
  \centering
  \includegraphics[width=\linewidth]{Filtre.png}
  \caption{Paramètres S simulés (HFSS) et mesurés d’un filtre 6 pôles et 4 zéros de transmission.}
  \label{fig:simu}
\end{figure}
Portez une attention particulière à la lisibilité des courbes, abscisses et ordonnées sur chacune de vos figures (cf. Figure~\ref{fig:simu}). Pour un meilleur rendu et un poids réduit, utilisez des figures vectorielles (pdf ou eps).

Il va de soi que vous pouvez mettre autant de paragraphes que vous le souhaitez, et qu'il est judicieux d'étayer votre propos par des figures claires et bien choisies.

\subsection{Tableaux}
Comme pour les figures, il est généralement conseillé de laisser Latex gérer le positionnement des tableaux pour éviter les blancs à la fin des colonnes comme c'est le cas pour le tableau~\ref{tab:marges}.
\begin{table}[ht]
	\centering
	\begin{tabular}{|m{0.3\linewidth}|m{0.2\linewidth}|}
	\hline
	\multicolumn{2}{|c|}{Marges}  \tabularnewline 
	\hline
	Haut & 2,5 cm \tabularnewline 
	\hline
	Bas & 2 cm \tabularnewline 
	\hline
	Gauche et droite & 2 cm \tabularnewline 
	\hline
	En-tête et pied de page & 1,25 cm \tabularnewline 
	\hline
	\end{tabular}
	\caption{Marges du document (format A4)}
	\label{tab:marges}
\end{table}

\subsection{Équations}
Les équations doivent être numérotées sur la droite de la façon suivante :
\begin{equation}
	y=ax+b
\end{equation}

\subsection{Références}
Les références seront rappelées dans le texte en utilisant des crochets \cite{kerherve, deplus}. La liste des références est détaillée en fin de texte selon les exemples fournis dans ce document.



%---[3. Un peu de culture]--------------- 
\section{Un peu de culture}
Le 24 octobre 1968, l'Université de Limoges est officiellement créée.

L’Etat a non seulement dialogué avec ses chercheurs, mais aussi avec le personnel politique local – Jacques Chirac, alors secrétaire d’Etat et député de la 3e circonscription de Corrèze, a appuyé cette naissance, comme Louis Longequeue, maire de la Ville. Le mouvement étudiant de Mai a également souligné l’urgence de créer de nouvelles universités dans le Pays. A Limoges  \cite{limoges}, il a permis d’exprimer le désir de la jeunesse de s’émanciper de la tutelle de l’Université de Poitiers.

De nouveaux locaux sortent de terre, comme l’IUT, construit en quelques mois seulement. A la rentrée 1968, l’Université de Limoges compte 4000 étudiants, dont 1266 en Droit, 939 en Médecine, 760 en Lettres et 132 à l’IUT.

Pendant une dizaine d’années, le développement des composantes va transformer radicalement Limoges, mais aussi la physionomie des villes accueillant tôt des sites distants, comme Egletons (1969).

Le 8 juillet 1968, le Collège des Sciences devient la Faculté des Sciences. Elle est aujourd’hui présente sur quatre sites. Le campus La Borie à Limoges avait été identifié par le maire de Limoges, Louis Longequeue – dit-on, par avion – comme idéal pour implanter les activités scientifiques.

Il propose aujourd’hui trois grands domaines de formations : Mathématiques-Informatique-Physique-Chimie ; Sciences de la vie et Sciences \& Techniques des Activités Physiques et Sportives (STAPS). Le site de Guéret est spécialisé en domotique. Le site d’Egletons est à la pointe en matière de génie civil, celui de Brive héberge notamment les STAPS, avec une spécialité pleine nature.

La FST compte aujourd’hui 3385 étudiants et rayonne par l’excellence des recherches menées par les 3 instituts XLIM, IPAM et GEIST auxquels elle est associée. Dès les années 1980, elle se distingue notamment des autres composantes par l’internationalisation avancée de ses étudiants – ils sont alors en proportion deux fois plus nombreux que dans le reste de l’Université. Ce sont alors le Maroc, la Tunisie, la Grande-Bretagne et le Liban qui fournissent les contingents les plus importants.



%---[4. Conclusion]--------------- 
\section{Conclusion}
Nous vous rappelons \textbf{que vous disposez de deux pages pour la 1ère soumission et de quatre pages pour la version finale informatique}. Libre à vous, après acceptation, d'étoffer les deux pages de soumission en donnant plus d'informations, c'est-à-dire en consolidant l'article par des résultats complémentaires \cite{deplus}, sans toutefois dépasser cette limite stricte de quatre pages. \textbf{La taille de 2 Mo pour le fichier PDF soumis devra être également respectée} afin que les actes de la conférence puissent être copiés sur un support numérique unique. En cas de non respect d'un de ces deux paramètres (nombre de pages ou taille de fichier), nous serions obligés d'utiliser la version courte pour publication dans les actes.



%---[5. Remerciements]--------------- 
\section{Remerciements (cas échéant)}
Nous remercions nos prédécesseurs de Bordeaux, Saint-Malo et Caen pour leur modèle et espérons que vous passerez de très bonnes JNM en notre compagnie à Limoges.



%\vfill

\vspace{0.5cm}
\bibliographystyle{IEEEtran}
\bibliography{biblio.bib}

\end{document}


\end{document}